\chapter{Experimental Setup}


\section{Hardware}

The experiments were conducted on a server equipped with an \textbf{AMD Ryzen Threadripper 3960X 24-Core Processor} clocked at \textbf{3.80 GHz}.
The processor features \textbf{24 cores}, each with \textbf{2 threads}, and is supported by \textbf{125,GiB of RAM}.
Cache specifications include \textbf{768,KiB L1d}, \textbf{768,KiB L1i}, \textbf{12,MiB L2}, and \textbf{128,MiB L3}.
The system is additionally equipped with two \textbf{NVIDIA GeForce RTX 2080 Ti GPUs}, each with \textbf{11,GiB of VRAM}.
The operating system used was \textbf{Ubuntu 22.04 LTS}.

The original project by \textcite{borassi_kadabra_2019} was implemented in \textbf{C++} and we compiled using \textbf{CMake version 4.2.1} with \textbf{GNU 11.4.0}.
Our implementation of KADABRA and BRAVA-GNN was implemented and run with \textbf{Julia 1.12.5}.
%The experiments were performed on an \textbf{Apple MacBook Pro (14-inch, 2023 model)},
%powered by an \textbf{Apple M2 Max System-on-a-Chip (SoC)}.
%The M2 Max features a \textbf{12-core CPU} (comprising 8 performance cores and 4 efficiency cores)
%and includes a \textbf{38-core GPU}, along with a \textbf{16-core Neural Engine}.
%The system was configured with \textbf{64\,GiB of unified memory} and a \textbf{1\,TB solid-state drive (SSD)}.
%The operating system was \textbf{macOS Sequoia 15.5}.
%The project was implemented in \textbf{C++20} and compiled using \textbf{CMake version 3.26.3} with \textbf{Apple Clang 17.0.0}.
\textbf{Simexpal 1.1}~\parencite{simexpal2019} was used to conduct the experiments.


\section{Implementation}\label{sec:implementation}
As basis for the implementation serves the framework for parametric maximum flow created by \textcite{beines_simpler_2024} in which we included the implementation of
pseudoflow and PPF by \textcite{Sauer_PPF_2025}.
Using these implementation of push relabel, pseudoflow, PBFS and PPF we implemented the Algorithms for DSP as described above.
For Greedy++ we based it on the implementation by \textcite{boob_flowless_2020}.
For a fairer comparison we adapted this implementation to use the same C++ datastructures and encapsulate it in a class, as
we did with the flow based algorithms.
Since constructing the flow network (Section~\ref{sec:reduction_to_flow}) plays an important role in the runtime for parametric maximum flow,
we therefore must also include the initiation of the necessary data structures for Greedy++.
%All algorithms receive the previously read edge list(see \ref{sec:Instances} for details).


\section[Instances]{Instances}\label{sec:Instances}

For comparability, we use the same instances as \textcite{noauthor_degree-mass_2026}.
A variety of instances are obtained from the SNAP database~\parencite{leskovec_snap_2014},
ASU’s Social Computing Data Repository~\parencite{zafarani_social_2009}, and the Koblenz Network Collection
\parencite{kunegis_konect_2013}.

All unweighted instances used in our experiments are listed in Table~\ref{tab:wei}
\begin{table}[h]
    \centering
    \resizebox{\textwidth}{!}{
        \begin{tabular}{lrrr}
\toprule
instance & $n$ & $m$ & $m/n$ \\
\midrule
p2p-Gnutella30 & 36\,682 & 88\,328 & 2.41 \\
email-Enron & 36\,692 & 367\,662 & 10.02 \\
musae-github & 37\,700 & 289\,003 & 7.67 \\
gemsec-Facebook & 50\,515 & 819\,306 & 16.22 \\
p2p-Gnutella31 & 62\,586 & 147\,892 & 2.36 \\
soc-Epinions1 & 75\,879 & 508\,837 & 6.71 \\
soc-Slashdot0811 & 77\,360 & 905\,468 & 11.70 \\
soc-Slashdot0902 & 82\,168 & 948\,464 & 11.54 \\
twitch-gamers & 168\,114 & 6\,797\,557 & 40.43 \\
email-EuAll & 265\,214 & 420\,045 & 1.58 \\
com-DBLP & 317\,080 & 1\,049\,866 & 3.31 \\
web-NotreDame & 325\,729 & 1\,497\,134 & 4.60 \\
web-BerkStan & 685\,230 & 7\,600\,595 & 11.09 \\
web-Google & 875\,713 & 5\,105\,039 & 5.83 \\
com-youtube & 1\,134\,890 & 2\,987\,624 & 2.63 \\
soc-Pokec & 1\,632\,803 & 30\,622\,564 & 18.75 \\
wiki-topcats & 1\,791\,489 & 28\,511\,807 & 15.92 \\
amazon & 2\,146\,057 & 5\,743\,146 & 2.68 \\
wiki-Talk & 2\,394\,385 & 5\,021\,410 & 2.10 \\
com-Orkut & 3\,072\,441 & 117\,185\,083 & 38.14 \\
cit-Patents & 3\,764\,117 & 16\,511\,741 & 4.39 \\
com-lj & 3\,997\,962 & 34\,681\,189 & 8.67 \\
dblp & 4\,000\,148 & 8\,649\,011 & 2.16 \\
soc-LiveJournal1 & 4\,847\,571 & 68\,993\,773 & 14.23 \\
\bottomrule
\end{tabular}

    }
    \caption{Overview of instances: $n$ is the number of vertices, $m$ the number of edges, 
    Sources: [1] KONECT~\parencite{kunegis_konect_2013}, [2] Zafrani~\parencite{zafarani_social_2009}; unmarked instances are from SNAP~\parencite{leskovec_snap_2014}.    }
    \label{tab:instance_table}
\end{table}

\section[Algorithm Performance Metrics]{Algorithm Performance Metrics}
\label{sec:algorithm-performance-metrics}
%Additionally, for  \emph{IPC} and \emph{DSPwithIntersect} algorithms, we recorded the number of iterations performed,
We conducted a detailed performance analysis of each exact algorithm on every graph instance.
For the exact methods, we recorded two primary metrics:
\begin{enumerate}
    \item the time required to construct the flow network,
    \item the runtime of the algorithm.
\end{enumerate}

For \emph{Greedy++}, we evaluated its performance as an approximation algorithm.
Given the optimal solution value, we specified target percentages
$(0.80, 0.90, 0.95, 0.99, 1.00)$ and measured both runtime and number of iterations
until the algorithm achieved at least each target fraction of the optimum.
Following \textcite{boob_flowless_2020} and confirmed by our experiments that
three iterations always reach at least 90\% of the optimum for the unweighted instances,
we also executed Greedy++ with exactly three iterations and recorded its runtime.

All experiments were repeated on both the original and the \emph{pruned} graphs
(see Chapter~\ref{ch:pruning}).
For pruning we tracked the pruning time and the resulting graph size.

Each algorithm–instance pair was run five times with a 30-minute time limit;
we report the median of these runs.
For pruning statistics we used the median of all instances and runs.

Throughout this thesis we use the following shorthand:
\begin{itemize}
    \item \textbf{PF}: IPC or {\DSPwithR} using pseudoflow, depending on context.
    \item \textbf{PR}: IPC or {\DSPwithR} using push–relabel, depending on context.
    \item \textbf{FPPF}: {\DSPwithFBA} using parametric pseudoflow.
    \item \textbf{FPBFS}: {\DSPwithFBA} using parametric breadth-first search.
    \item \textbf{G-XX}: Greedy++ run until XX\% of the optimum is reached (e.g., G-90).
    \item \textbf{G-3}: Greedy++ run for exactly three iterations.
    \item \textbf{P-X}: Pruning Version.
\end{itemize}





%
%livehournal groupmemberships timeout bei peeler100 und pbfs->immernoch, aber das einzige
%orkut-links bei peeler100 -> weiterhin.
%twitter follows bei fpbfs-> immernoch, aber nur bei nicht reduced.
%
%Hab jetzt nur die instanzen aus flowless: also alle entfernt, besonders diese hier, weil funtkionierte nicht:
%Instanz rausgenommen: US-Airport_2010 bei Peeler 80 weighted
%US-Airport500 komische density 12184817 bei Peeler 80 weighted
